% -------------------------------
% ТЕРМИНЫ И ОПРЕДЕЛЕНИЯ
% -------------------------------

\newglossaryentry{id1}{
    name={\LaTeX},
    description={система компьютерной вёрстки, используемая для подготовки научных, технических и учебных документов с поддержкой формул, таблиц, ссылок и автоматической структуры документа}
}

\newglossaryentry{id2}{
    name={МАИ},
    description={Московский авиационный институт, федеральное государственное автономное образовательное учреждение высшего образования}
}

\newglossaryentry{id3}{
    name={$S^{2}$},
    description={математическое обозначение квадрата величины $S$; используется в формулах и перечне обозначений при описании математических зависимостей}
}

\newglossaryentry{id4}{
    name={SQL},
    description={Structured Query Language, язык структурированных запросов, применяемый для создания, изменения, выборки и анализа данных в реляционных базах данных}
}

\newglossaryentry{id5}{
    name={Реляционная база данных},
    description={база данных, в которой информация представлена в виде связанных таблиц, состоящих из строк и столбцов, а доступ к данным осуществляется с помощью SQL-запросов}
}

\newglossaryentry{id6}{
    name={Естественный язык},
    description={язык повседневного общения пользователя, например русский или английский, используемый для формулирования запроса без необходимости знания SQL}
}

\newglossaryentry{id7}{
    name={Интерфейс естественного языка к базам данных},
    description={программный интерфейс, позволяющий пользователю обращаться к базе данных с помощью текстовых вопросов на естественном языке вместо ручного написания SQL-запросов}
}

\newglossaryentry{id8}{
    name={Text-to-SQL},
    description={задача автоматического преобразования пользовательского запроса, сформулированного на естественном языке, в исполняемый SQL-запрос}
}

\newglossaryentry{id9}{
    name={Интеллектуальный модуль интерактивного анализа данных},
    description={разрабатываемое программное решение, обеспечивающее взаимодействие пользователя с реляционной базой данных на естественном языке с применением больших языковых моделей, RAG-архитектуры и механизма самокоррекции}
}

\newglossaryentry{id10}{
    name={Большая языковая модель},
    description={нейросетевая модель, обученная на больших объёмах текстовых и программных данных и способная выполнять задачи понимания естественного языка, генерации текста и создания программного кода, включая SQL-запросы}
}

\newglossaryentry{id11}{
    name={RAG},
    description={подход к генерации ответа, при котором языковая модель предварительно получает релевантный контекст, извлечённый из базы знаний или векторного хранилища}
}

\newglossaryentry{id12}{
    name={RAG-конвейер},
    description={последовательность операций, включающая векторизацию запроса пользователя, поиск релевантного контекста и передачу найденных фрагментов в языковую модель для генерации ответа или SQL-запроса}
}

\newglossaryentry{id13}{
    name={Конвейер индексации},
    description={последовательность операций предварительной обработки данных, включающая загрузку документа, извлечение текста, разбиение на фрагменты, векторизацию и сохранение результатов в векторной базе данных}
}

\newglossaryentry{id14}{
    name={Конвейер поиска},
    description={последовательность операций обработки пользовательского запроса, включающая векторизацию запроса, поиск ближайших по смыслу фрагментов и передачу найденного контекста в языковую модель}
}

\newglossaryentry{id15}{
    name={База знаний},
    description={совокупность проиндексированных документов, описаний схемы базы данных, метаинформации, примеров данных и иных текстовых фрагментов, используемых системой для поиска релевантного контекста}
}

\newglossaryentry{id16}{
    name={Метаинформация схемы базы данных},
    description={дополнительные сведения о структуре базы данных, включающие названия таблиц и столбцов, типы данных, комментарии, связи между таблицами и репрезентативные примеры значений}
}

\newglossaryentry{id17}{
    name={Схема базы данных},
    description={описание структуры базы данных, включающее таблицы, столбцы, типы данных, ключи, ограничения и связи между сущностями}
}

\newglossaryentry{id18}{
    name={Связывание схемы},
    description={процесс сопоставления пользовательского вопроса с релевантными таблицами, столбцами и связями базы данных, необходимыми для построения корректного SQL-запроса}
}

\newglossaryentry{id19}{
    name={Векторизация текста},
    description={процесс преобразования текстового фрагмента в числовой вектор фиксированной размерности, отражающий его смысловое содержание}
}

\newglossaryentry{id20}{
    name={Эмбеддинг},
    description={числовое векторное представление текста, позволяющее сравнивать фрагменты документов и запросы пользователя по семантической близости}
}

\newglossaryentry{id21}{
    name={Векторное хранилище},
    description={специализированная база данных для хранения эмбеддингов и выполнения поиска ближайших векторов по выбранной метрике сходства}
}

\newglossaryentry{id22}{
    name={Семантический поиск},
    description={метод поиска информации, основанный не только на совпадении ключевых слов, но и на смысловой близости между запросом пользователя и проиндексированными фрагментами текста}
}

\newglossaryentry{id23}{
    name={Косинусное сходство},
    description={метрика близости двух векторов, основанная на вычислении косинуса угла между ними и применяемая для оценки семантической близости текстовых фрагментов}
}

\newglossaryentry{id24}{
    name={Чанк},
    description={отдельный фрагмент текста, полученный после разбиения документа или описания схемы базы данных и используемый как минимальная единица индексации в RAG-конвейере}
}

\newglossaryentry{id25}{
    name={Разбиение текста на фрагменты},
    description={метод предварительной обработки текста, при котором большой документ делится на небольшие смысловые части для последующей векторизации и более точного поиска}
}

\newglossaryentry{id26}{
    name={Агентный модуль},
    description={программный компонент интеллектуальной системы, использующий языковую модель и внешние инструменты для многошаговой обработки запроса, генерации SQL, проверки результата и принятия решения о самокоррекции}
}

\newglossaryentry{id27}{
    name={ReAct-паттерн},
    description={подход к построению интеллектуальных агентов, основанный на чередовании этапов рассуждения и выполнения действий с использованием внешних инструментов}
}

\newglossaryentry{id28}{
    name={Самокоррекция SQL-запроса},
    description={механизм автоматического исправления сгенерированного SQL-запроса на основе анализа ошибки выполнения, полученной от системы управления базами данных}
}

\newglossaryentry{id29}{
    name={Self-Heal},
    description={механизм самовосстановления или самокоррекции, при котором агент анализирует ошибку выполнения SQL-запроса и формирует исправленную версию запроса без участия пользователя}
}

\newglossaryentry{id30}{
    name={Обратная связь от СУБД},
    description={диагностическая информация, возвращаемая системой управления базами данных при выполнении SQL-запроса, включая сообщения об ошибках синтаксиса, отсутствующих таблицах или неверных именах полей}
}

\newglossaryentry{id31}{
    name={Галлюцинация языковой модели},
    description={ошибка генеративной модели, при которой она формирует правдоподобный, но фактически неверный ответ, например использует несуществующие таблицы, поля или связи базы данных}
}

\newglossaryentry{id32}{
    name={Микросервисная архитектура},
    description={архитектурный подход, при котором система разделяется на независимые сервисы, каждый из которых выполняет отдельную функцию и взаимодействует с другими сервисами через программные интерфейсы}
}

\newglossaryentry{id33}{
    name={Объектное хранилище},
    description={тип хранилища данных, предназначенный для сохранения файлов как отдельных объектов с метаданными; в работе используется для хранения исходных документов}
}

\newglossaryentry{id34}{
    name={Промпт},
    description={текстовая инструкция или входной запрос, передаваемый в языковую модель для управления процессом генерации ответа или SQL-запроса}
}

\newglossaryentry{id35}{
    name={SQL-запрос},
    description={формализованная инструкция на языке SQL, предназначенная для получения, фильтрации, агрегации, соединения или анализа данных в реляционной базе данных}
}

\newglossaryentry{id36}{
    name={Тестовая SQL-задача},
    description={контрольный пользовательский запрос с эталонным SQL-решением, применяемый для оценки качества генерации SQL-запросов языковыми моделями}
}

\newglossaryentry{id37}{
    name={Метрика качества},
    description={количественный показатель, используемый для оценки корректности, точности или полезности сгенерированного SQL-запроса}
}

